%% LyX 2.3.6.2 created this file.  For more info, see http://www.lyx.org/.
%% Do not edit unless you really know what you are doing.
\documentclass[english]{article}
\usepackage[T1]{fontenc}
\usepackage[latin9]{inputenc}
\usepackage{amsmath}
\usepackage{amssymb}
\usepackage{babel}
\begin{document}
As detailed in the main text, the starting point of the resistivity
calculation is the variational Boltzmann approach justified by the
long mean free path $k_{F}\ell$. Within this approach the transport
lifetime can estimated by 
\begin{equation}
\rho\leq\frac{\frac{V}{2kT}\sum_{\boldsymbol{q}\boldsymbol{p}}P_{\boldsymbol{p}\boldsymbol{q}}\left(\psi_{\boldsymbol{q}}-\psi_{\boldsymbol{p}}\right)^{2}}{\left|\sum_{\boldsymbol{q}}e\frac{\partial f}{\partial\epsilon_{\boldsymbol{q}}}\psi_{\boldsymbol{q}}\boldsymbol{v}\left(\boldsymbol{q}\right)\cdot\boldsymbol{n}\right|^{2}}
\end{equation}

with $\psi_{\boldsymbol{q}}$ a variational anzatz, $P_{\boldsymbol{p}\boldsymbol{q}}$
the transition probaility from different momentum states,$V$ the
volume, $\boldsymbol{n}$ is the direction of the electric field $\boldsymbol{E}$,
$f$ the fermi-dirac distribution, $e$ the electron charge , $\boldsymbol{v}\left(\boldsymbol{q}\right)=\partial_{\boldsymbol{q}}\epsilon_{\boldsymbol{q}},$
and $\epsilon_{\boldsymbol{q}}$ the dispersion. Using the ansatz
$\psi_{\boldsymbol{q}}=-e\boldsymbol{v}_{\boldsymbol{q}}\cdot\boldsymbol{E}\tau_{\text{tr}}$
we recover the drude formula with the momentum relaxation time
\[
\tau_{\text{tr}}^{-1}=\frac{\frac{1}{kT}\sum_{\boldsymbol{q}\boldsymbol{p}}P_{\boldsymbol{p}\boldsymbol{q}}\left(\boldsymbol{v}\left(\boldsymbol{q}\right)\cdot\boldsymbol{n}-\boldsymbol{v}\left(\boldsymbol{q}\right)\cdot\boldsymbol{n}\right)^{2}}{\sum_{\boldsymbol{q}}2\left(-\frac{\partial f}{\partial\epsilon_{\boldsymbol{q}}}\right)\left(\boldsymbol{v}\left(\boldsymbol{q}\right)\cdot\boldsymbol{n}\right)^{2}}.
\]

Assuming Mathiessen's rule holds, we have 
\begin{equation}
\tau_{\text{tr}}^{-1}=\tau_{\text{tr},K}^{-1}+\tau_{\text{tr},EME}^{-1}+\tau_{\text{tr},ep}^{-1}
\end{equation}

Where each individual scattering rate is determined by the transition
probabilities 
\begin{equation}
P_{\boldsymbol{p}\boldsymbol{q}}^{K}=kT\frac{2\pi}{\hbar}\left(-\frac{\partial f_{\boldsymbol{p}}}{\partial\epsilon_{\boldsymbol{p}}}\right)\left|J_{K}\left(\boldsymbol{p},\boldsymbol{q}\right)\right|^{2}\chi''\left(\boldsymbol{p}-\boldsymbol{q},\varepsilon_{\boldsymbol{p}}-\epsilon_{\boldsymbol{q}}\right)\left(n_{B}\left(\varepsilon_{\boldsymbol{p}}-\epsilon_{\boldsymbol{q}}\right)+n_{F}\left(-\epsilon_{\boldsymbol{q}}\right)\right)
\end{equation}

\begin{equation}
P_{\boldsymbol{k}\boldsymbol{p}}^{EME}=kT\left(-\frac{\partial f}{\partial\epsilon_{\boldsymbol{p}}}\right)\left|\tilde{\alpha}\left(\boldsymbol{p},\boldsymbol{q}\right)\right|^{2}\sum_{\boldsymbol{k}}\int\frac{d\nu}{\pi}\chi''\left(\boldsymbol{k}+\boldsymbol{p}-\boldsymbol{q},\epsilon_{\boldsymbol{p}}-\nu\right)\mathcal{B}\left(\boldsymbol{k},\nu-\varepsilon_{\boldsymbol{q}}\right)\left[n_{B}\left(\epsilon_{\boldsymbol{p}}-\nu\right)+n_{F}\left(-\nu\right)\right]\left[n_{B}\left(\nu-\varepsilon_{\boldsymbol{q}}\right)+n_{F}\left(-\varepsilon_{\boldsymbol{q}}\right)\right]
\end{equation}

\begin{equation}
P_{\boldsymbol{p}\boldsymbol{q}}^{ep}=kT\frac{2\pi}{\hbar}\left(-\frac{\partial f_{\boldsymbol{p}}}{\partial\epsilon_{\boldsymbol{p}}}\right)\sum_{\lambda}\left|\alpha_{\lambda}\left(\boldsymbol{p},\boldsymbol{q}\right)\right|^{2}\mathcal{B}_{\lambda}\left(\boldsymbol{p}-\boldsymbol{q},\varepsilon_{\boldsymbol{p}}-\epsilon_{\boldsymbol{q}}\right)\left(n_{B}\left(\varepsilon_{\boldsymbol{p}}-\epsilon_{\boldsymbol{q}}\right)+f\left(-\epsilon_{\boldsymbol{q}}\right)\right)
\end{equation}

where $\chi''\left(\boldsymbol{p},\omega\right)$ is the spin susceptibility
and $\mathcal{B}_{\lambda}\left(\boldsymbol{k},\nu\right)$ is the
phonon spectral function for mode $\lambda=o,a$ which can be either
acoustic or optical. We also reintroduced the momentum dependence
in all the couplings. In the case where the magnetic and itinerant
electron layers are perfectly aligned we have 
\begin{equation}
J_{K}\left(\boldsymbol{p},\boldsymbol{q}\right)=J_{K}\left[2\cos\left(p_{z}a_{z}\right)\right]\left[2\cos\left(q_{z}a_{z}\right)\right],
\end{equation}

\begin{equation}
\tilde{\alpha}\left(\boldsymbol{p},\boldsymbol{q}\right)=\tilde{\alpha}\left[2\cos\left(p_{z}a_{z}\right)\right]\left[2\cos\left(q_{z}a_{z}\right)\right],
\end{equation}

\begin{equation}
\alpha_{o}\left(\boldsymbol{p},\boldsymbol{q}\right)=\alpha_{o},
\end{equation}

and 

\begin{equation}
\alpha_{a}\left(\boldsymbol{p},\boldsymbol{q}\right)=\alpha_{a}\left|\boldsymbol{p}-\boldsymbol{q}\right|.
\end{equation}

Where the coupling constants where defined in the main text. Additonally
we consider the local approximation for 
\begin{equation}
\chi''\left(\nu\right)=S\left(S+1\right)\beta\nu\frac{\theta\left(2\pi J_{H}-\left|\nu\right|\right)}{2J_{H}}
\end{equation}

with the sum rule constraint 
\begin{equation}
\int\frac{\chi''\left(\nu\right)}{\beta\nu}d\nu=2\pi S\left(S+1\right).
\end{equation}

We consider the spectral function for free phonons given by 
\begin{equation}
\mathcal{B}_{\lambda}\left(\boldsymbol{k},\nu\right)=\frac{\pi\hbar}{2M\omega_{\boldsymbol{k},\lambda}}\left(\delta\left(\nu-\omega_{\boldsymbol{k},\lambda}\right)-\delta\left(\nu+\omega_{\boldsymbol{k},\lambda}\right)\right)
\end{equation}

where for optical phonons we have $\omega_{\boldsymbol{k},o}=\omega_{0}$
while for acoustic phonons $\omega_{\boldsymbol{k},a}=ck$. The integrals
above can then be evaluated analytically if we take a quasi-2D circular
fermi surface and the scattering rate becomes
\begin{align}
\tau_{tr}^{-1} & =2\pi\frac{\lambda_{tr,EP,o}}{\hbar\beta}\mathcal{I}_{1}\left(\frac{\beta\hbar\omega_{0}}{2}\right)+2\pi\frac{\lambda_{tr,EP,a}}{\hbar\beta}\mathcal{I}_{2}\left(\frac{\beta\hbar\omega_{BG}}{2}\right)\\
 & +2\pi\frac{\lambda_{K}J_{K}}{\hbar}\mathcal{I}_{3}\left(2\pi\beta J_{H}\right)+2\pi\frac{\lambda_{tr,EME}}{\hbar\beta}\mathcal{I}_{4}\left(2\pi\beta J_{H},\beta\hbar\tilde{\omega}_{0}\right)
\end{align}

Explicitly we have the definitions:

\begin{equation}
\mathcal{I}_{1}\left(x\right)=\left[x\text{csch}\left(x\right)\right]^{2},
\end{equation}

which is the usual result for optical phonons,

\begin{equation}
\mathcal{I}_{2}\left(x\right)=\frac{4}{\pi}\frac{1}{x^{3}}\int_{0}^{x}dy\ \frac{y^{4}}{\sqrt{1-\left(\frac{y}{x}\right)^{2}}}\left[\text{csch}\left(y\right)\right]^{2},
\end{equation}

which is the result for acousic phonons interacting with a circular
Fermi surface in two dimensions,

\begin{align}
\mathcal{I}_{3}\left(y\right) & =-\frac{4\text{Li}_{2}\left(e^{-y}\right)}{y}-\frac{2y}{e^{y}-1}-4y+\frac{2\pi^{2}}{3y}+4\log\left(e^{y}-1\right),
\end{align}

corresponds to the pure kondo scattering at high temperatures, whre
the logarithmic integrals result from integrating the fermi function,
and finally 

\begin{equation}
\mathcal{I}_{4}\left(y,z\right)=\frac{z^{2}}{\cosh\left(z\right)-1}\mathcal{J}\left(y,z\right),
\end{equation}

is the result of the integrals involving the EME coupling, with the
additional definition of:

\begin{align}
\mathcal{J}\left(y,z\right) & =\frac{-2z\text{Li}_{2}\left(e^{-y}\right)+\left(2y+z\right)\text{Li}_{2}\left(e^{-y-z}\right)+\left(z-2y\right)\text{Li}_{2}\left(e^{z-y}\right)+2\text{Li}_{3}\left(e^{-y-z}\right)-2\text{Li}_{3}\left(e^{z-y}\right)}{2yz}\\
 & +\frac{z^{2}}{12y}-\frac{1}{2}\log\left(\left(e^{y}-e^{z}\right)\left(e^{y+z}-1\right)\right)+\frac{y\log\left(\sinh\left(\frac{y-z}{2}\right)\text{csch}\left(\frac{y+z}{2}\right)\right)}{2z}+\frac{y}{2}+\frac{\pi^{2}}{3y}+\log\left(e^{y}-1\right)+\frac{z}{2}.
\end{align}

These functions have the following high temperature limits which can
be obtained from the asymptotic expansions for the polylogarithm
\begin{equation}
\lim_{T\rightarrow\infty}\mathcal{I}_{1}\left(\frac{\beta\hbar\omega_{0}}{2}\right)=\lim_{T\rightarrow\infty}\mathcal{I}_{2}\left(\frac{\beta\hbar\omega_{BG}}{2}\right)=1
\end{equation}

and 
\begin{equation}
\lim_{T\rightarrow\infty}\mathcal{I}_{3}\left(2\pi\beta J_{H}\right)=\lim_{T\rightarrow\infty}\mathcal{I}_{4}\left(2\pi\beta J_{H},\beta\hbar\tilde{\omega}_{0}\right)=2.
\end{equation}

In the main text, we use the effective mass of PdCoO$_{2}$ extracted
from quantum oscillations given in {[}THESIS{]}, and the density corresponding
to half filling to extract the transport rate from experiments following
the procedure in {[}ANDY anlysis{]}. Then we fit the resulting scattering
rate to the first two terms which correspond to electron-phonon scattering
in $\tau_{tr}^{-1}$. With this procedure, we find the parameters
$\lambda_{tr,EP,o}=0.024$,$\hbar\omega_{0}=120$meV,$\lambda_{tr,EP,a}=0.043$,
and $\hbar\omega_{BG}=29$meV, which are in close agreement with the
reported values in {[}XX{]}. The difference between our treatment
and that in {[}XX{]} is that we consider a 2D fermi surface instead
of a 3D Fermi surface as in those works. At intermediate temperatures
where our theory is valid we don't see significant differences in
the fit quality between the model considered here and the ones in
{[}XX{]}.

We then carry out the same procedure for PdCrO$_{2}$ and substract
the resulting scattering rate from the scattering rate extracted from
PdCoO$_{2}$. This would isolate the last two terms that correspond
to Kondo and EME scattering in $\tau_{tr}^{-1}$. By fitting this
excess scattering rate we find the EME parameters $\lambda_{tr,EME}=0.025$
and $\hbar\tilde{\omega}_{0}=45$meV. Since we use the high temperature
approximation for $\chi''$ we don't fit the low temperature regime
of $\Delta_{\tau}=\tau_{tr,Cr}^{-1}-\tau_{tr,Co}^{-1}$ but only carry
the fit at temperatures above 80K ($J_{H}\sim72K$). Note that the
constant offset between the PdCoO$_{2}$ and the PdCrO$_{2}$ resistivities
is not fitted. This excess scattering rate comes purely from the Kondo
scattering with the dimensionless coupling $\lambda_{K}=4S\left(S+1\right)J_{K}\nu_{0}$
and a saturation scale set by $J_{H}$, where the value of these couplings
is listed above. 

In the above, one can trace the variuous features of $\rho$ for PdCrO$_{2}$
to each of the terms in {[}TAUEQ{]}.The Kondo term increases sharply
and completely saturates at a temperature scale set by $2\pi J_{H}$.
This approach to saturation would be responsible for the sublinear
approach to the $T$-linear regime. This sublinear piece dominates
at $T\sim J_{H}$ over the contributions from the phonon and the magnetoelastic
scattering rates. Eventually, both of these contributions catch up
with the Kondo scattering rate and the resistivity becomes $T$ linear
which can be seen from the high temperature limits of $\mathcal{I}_{4}$
and $\mathcal{I}_{3}$ and the expression for $\tau_{tr}^{-1}$. Taking
the high temperature limit, we can determine the coefficient of the
linear in $T$ scattering rate. Using the asymptotic expressions above
and the values of the fitted dimensionless couplings we get 
\begin{align}
\lim_{T\rightarrow\infty}\frac{d}{dT}\tau_{tr}^{-1} & =\left(\frac{k_{B}}{\hbar}\right)2\pi\left(\lambda_{tr,EP,o}+\lambda_{tr,EP,a}+2\lambda_{tr,EME}\right)\\
 & \sim0.75\left(\frac{k_{B}}{\hbar}\right)
\end{align}

which is close to the Planckian value. We interpret this as a coincidence
as the scattering is quasielastic in this high temperature regime. 

Finally, we note that neither $\tau_{ep}^{-1}$ nor $\tau_{EME}^{-1}$
are in the fully saturated regime for the values of the fit parameters
obtained from the above procedure (see Fig.XX). In particular. Note
that, by itself, PdCoO$_{2}$ show a non-linear resistivity with a
high saturation temperature given the high frequency of the in plane
optical mode. At the same time the excess resistivity from the EME
interaction is generated by a mode that also saturates at relatively
large temperatures. The addition of the pure Kondo scattering gives
an additional sublinear contribution to the excess resistivity. In
the end, the superlinear scattering rate from electron phonon scattering
and the sublinear contribution from the EME+Kondo interaction balance
out to give an almost linear resistivity for PdCrO$_{2}.$
\end{document}
